\documentclass[12pt, a4paper]{article}
% ============================================================
%  preamble.tex — تنظیمات مشترک اسناد LaTeX (کامپایل با XeLaTeX)
% ============================================================
\usepackage{xcolor}
\usepackage{geometry}
\geometry{margin=2.5cm}
\usepackage{titlesec}
\usepackage{enumitem}
\setlist{nosep, leftmargin=*}
\usepackage{longtable}
\usepackage{booktabs}
\usepackage{array}
\usepackage{calc}
\usepackage{ragged2e}
\usepackage{fancyhdr}
\usepackage{etoolbox}
\usepackage{footnote}

% ---------- لازمهٔ جدول‌های خروجی pandoc ----------
\makeatletter
\patchcmd\longtable{\par}{\if@noskipsec\mbox{}\fi\par}{}{}
\makeatother
\makesavenoteenv{longtable}
\providecommand{\tightlist}{\setlength{\itemsep}{0pt}\setlength{\parskip}{0pt}}

% ---------- ظاهر ----------
\definecolor{primary}{RGB}{31, 78, 121}
\definecolor{accent}{RGB}{0, 150, 136}
\titleformat{\section}{\Large\bfseries\color{primary}}{\thesection}{1em}{}
\titleformat{\subsection}{\large\bfseries}{\thesubsection}{1em}{}
\titleformat{\subsubsection}{\normalsize\bfseries}{\thesubsubsection}{1em}{}
\renewcommand{\arraystretch}{1.3}
\linespread{1.3}

\pagestyle{fancy}
\fancyhf{}
\fancyfoot[C]{\thepage}
\renewcommand{\headrulewidth}{0.4pt}
\setlength{\headheight}{14pt}

% ---------- hyperref قبل از xepersian ----------
\usepackage{hyperref}
\hypersetup{colorlinks=true, linkcolor=primary, urlcolor=primary, citecolor=primary}

% ---------- فارسی‌سازی (آخرین بسته) ----------
\usepackage{xepersian}
\settextfont{Tahoma}
\setlatintextfont{Times New Roman}
\setdigitfont{Tahoma}

% ---------- شمارندهٔ لازم برای جدول‌های pandoc (LTcaptype = none) ----------
\newcounter{none}

\fancyhead[R]{\small داستان‌های کاربر — کلاس‌یار}
\fancyhead[L]{\small نسخهٔ ۱.۰}

\begin{document}

\begin{titlepage}
\centering
\vspace*{3cm}
{\Huge\bfseries داستان‌های کاربر\par}
\vspace{0.5cm}
{\LARGE\bfseries پروژهٔ کلاس‌یار\par}
\vspace{0.5cm}
{\large پلتفرم اشتراک جزوه، نمونه‌سؤال و ارزیابی استادان\par}
\vspace{2cm}
{\large تهیه‌کنندگان:\par}
{\large مریم قندی زنجانی، ملیکا شهبازی، آنیسا فروغمند، مائده رضایی، محمدمهدی ذوالقدر\par}
\vspace{0.5cm}
{\large دانشگاه آزاد زنجان\par}
\vspace{1cm}
{\large ۱۴۰۵/۰۷/۱۱\par}
\end{titlepage}

\tableofcontents
\newpage

{\def\LTcaptype{none} % do not increment counter
\begin{longtable}[]{@{}
  >{\raggedright\arraybackslash}p{(\linewidth - 2\tabcolsep) * \real{0.5000}}
  >{\raggedright\arraybackslash}p{(\linewidth - 2\tabcolsep) * \real{0.5000}}@{}}
\toprule\noalign{}
\endhead
\bottomrule\noalign{}
\endlastfoot
\textbf{نسخهٔ سند} & ۱.۰ \\
\textbf{وضعیت} & پیش‌نویس — ۱۰ Story اولیه \\
\textbf{مرجع} & \texttt{docs/01-vision.md} و \texttt{docs/02-requirements/personas.md} \\
\end{longtable}
}

\textbf{قالب:} به‌عنوان {[}نقش{]} می‌خواهم {[}هدف{]} تا {[}ارزش{]}.
\textbf{اولویت:} MoSCoW (Must / Should / Could). \textbf{ارتباط:} شمارهٔ قابلیت در سند چشم‌انداز.

{\def\LTcaptype{none} % do not increment counter
\begin{longtable}[]{@{}
  >{\raggedright\arraybackslash}p{(\linewidth - 10\tabcolsep) * \real{0.0690}}
  >{\raggedright\arraybackslash}p{(\linewidth - 10\tabcolsep) * \real{0.1149}}
  >{\raggedright\arraybackslash}p{(\linewidth - 10\tabcolsep) * \real{0.5747}}
  >{\raggedright\arraybackslash}p{(\linewidth - 10\tabcolsep) * \real{0.0805}}
  >{\raggedright\arraybackslash}p{(\linewidth - 10\tabcolsep) * \real{0.0690}}
  >{\raggedright\arraybackslash}p{(\linewidth - 10\tabcolsep) * \real{0.0920}}@{}}
\toprule\noalign{}
\begin{minipage}[b]{\linewidth}\raggedright
شناسه
\end{minipage} & \begin{minipage}[b]{\linewidth}\raggedright
نقش
\end{minipage} & \begin{minipage}[b]{\linewidth}\raggedright
داستان
\end{minipage} & \begin{minipage}[b]{\linewidth}\raggedright
اولویت
\end{minipage} & \begin{minipage}[b]{\linewidth}\raggedright
Sprint
\end{minipage} & \begin{minipage}[b]{\linewidth}\raggedright
قابلیت
\end{minipage} \\
\midrule\noalign{}
\endhead
\bottomrule\noalign{}
\endlastfoot
US-01 & دانشجو / استاد & به‌عنوان کاربر جدید می‌خواهم ثبت‌نام کنم و وارد شوم تا از امکانات مخصوص نقشم استفاده کنم & Must & 1 & F-01 \\
US-02 & امیرحسین & به‌عنوان دانشجوی مشارکت‌کننده می‌خواهم جزوه‌ام را با مشخصات استاد، درس، واحد و ترم آپلود کنم تا دیگران پیدایش کنند & Must & 1 & F-03 \\
US-03 & نگار & به‌عنوان دانشجو می‌خواهم جزوه‌ها را بر پایهٔ استاد، استان، واحد و درس فیلتر کنم تا سریع جزوهٔ موردنظرم را بیابم & Must & 1-2 & F-04 \\
US-04 & نگار & به‌عنوان دانشجو می‌خواهم جزئیات جزوه را ببینم و آن را دانلود کنم تا قبل از مطالعه از مناسب‌بودنش مطمئن شوم & Must & 1 & F-05 \\
US-05 & نگار & به‌عنوان دانشجو می‌خواهم به جزوه امتیاز بدهم و نظر بنویسم تا بقیه کیفیتش را بدانند & Must & 2 & F-06, F-09 \\
US-06 & نگار & به‌عنوان دانشجو می‌خواهم صفحهٔ هر استاد را ببینم و به او امتیاز بدهم تا دانشجویان بعدی او را بهتر بشناسند & Should & 2 & F-07, F-08 \\
US-07 & نگار & به‌عنوان دانشجو می‌خواهم نمونه‌سؤال‌ها و لینک گروه کلاس هر استاد را ببینم تا بهتر آماده شوم & Should & 3 & F-10 \\
US-08 & دکتر کاظمی & به‌عنوان استاد می‌خواهم جزوه‌های درسم را تأیید کنم و برای آپلودکننده نمرهٔ تشویقی ثبت کنم تا دانشجویان فعال تشویق شوند & Could & 3 & F-11 \\
US-09 & خانم نوروزی & به‌عنوان مدیر می‌خواهم گزارش‌های محتوای نامناسب را ببینم و آن را حذف یا رد کنم تا اعتبار سامانه حفظ شود & Must & 3 & F-12, F-13 \\
US-10 & خانم نوروزی & به‌عنوان مدیر می‌خواهم استان، دانشگاه، واحد، درس و استاد را تعریف و درخواست‌های افزودن را تأیید کنم تا فیلترها درست کار کنند & Must & 1 & F-02 \\
\end{longtable}
}

\begin{center}\rule{0.5\linewidth}{0.5pt}\end{center}

\section{معیارهای پذیرش (Acceptance Criteria)}

\subsection{US-01 — ثبت‌نام و ورود}

\begin{itemize}
\tightlist
\item
  کاربر با ایمیل، گذرواژه و انتخاب نقش (دانشجو یا استاد) ثبت‌نام می‌کند.
\item
  گذرواژه با bcrypt ذخیره می‌شود و گذرواژهٔ کمتر از ۸ کاراکتر پذیرفته نمی‌شود.
\item
  حساب استاد پس از تأیید مدیر فعال می‌شود.
\end{itemize}

\subsection{US-02 — آپلود جزوه}

\begin{itemize}
\tightlist
\item
  انتخاب استاد، درس، استان، واحد و ترم اجباری است و توضیح اختیاری.
\item
  فقط PDF، JPG، PNG و DOCX تا ۲۰ مگابایت پذیرفته می‌شود.
\item
  پس از آپلود، جزوه در فهرست آپلودکننده دیده می‌شود.
\end{itemize}

\subsection{US-03 — جست‌وجو و فیلتر}

\begin{itemize}
\tightlist
\item
  ترکیب چند فیلتر (استاد + استان + واحد) درست کار می‌کند.
\item
  نتیجه کمتر از ۲ ثانیه نمایش داده می‌شود.
\item
  اگر نتیجه‌ای نبود، پیام «جزوه‌ای یافت نشد» نمایش داده می‌شود.
\end{itemize}

\subsection{US-04 — جزئیات و دانلود}

\begin{itemize}
\tightlist
\item
  عنوان، استاد، درس، ترم، آپلودکننده، میانگین امتیاز و تعداد دانلود نمایش داده می‌شود.
\item
  دکمهٔ دانلود فایل اصلی را بدون خطا می‌دهد.
\end{itemize}

\subsection{US-05 — امتیاز و نظر}

\begin{itemize}
\tightlist
\item
  امتیاز بین ۱ تا ۵ است و هر کاربر برای هر جزوه فقط یک بار می‌تواند امتیاز بدهد.
\item
  میانگین امتیاز بلافاصله به‌روز می‌شود.
\item
  نظر خالی ثبت نمی‌شود و دیگران می‌توانند زیر آن پاسخ بدهند.
\end{itemize}

\subsection{US-06 — صفحهٔ استاد و امتیاز به استاد}

\begin{itemize}
\tightlist
\item
  صفحهٔ استاد جزوه‌ها، میانگین امتیاز و نظرها را نشان می‌دهد.
\item
  هر دانشجو برای هر استاد فقط یک امتیاز ثبت می‌کند و می‌تواند آن را ویرایش کند.
\end{itemize}

\subsection{US-07 — نمونه‌سؤال و لینک گروه}

\begin{itemize}
\tightlist
\item
  دانشجو می‌تواند نمونه‌سؤال (فایل) یا لینک گروه کلاس را برای یک استاد و درس ثبت کند.
\item
  لینک‌ها فقط با ساختار معتبر (https) ذخیره می‌شوند.
\end{itemize}

\subsection{US-08 — تأیید و نمرهٔ تشویقی}

\begin{itemize}
\tightlist
\item
  استاد فقط جزوه‌های مربوط به خودش را می‌بیند و تأیید می‌کند.
\item
  نمرهٔ تشویقی همراه یادداشت کوتاه ثبت می‌شود و آپلودکننده آن را در پروفایلش می‌بیند.
\item
  سامانه صریحاً نشان می‌دهد که این نمره رسمی نیست و فقط تشویقی است.
\end{itemize}

\subsection{US-09 — گزارش و بررسی}

\begin{itemize}
\tightlist
\item
  هر کاربر می‌تواند جزوه یا نظر را با ذکر دلیل گزارش کند.
\item
  مدیر فهرست گزارش‌ها را می‌بیند و محتوا را حذف یا گزارش را رد می‌کند.
\item
  وضعیت گزارش (جدید، بررسی‌شده، رد‌شده) ثبت می‌شود.
\end{itemize}

\subsection{US-10 — اطلاعات پایه}

\begin{itemize}
\tightlist
\item
  مدیر استان، دانشگاه، واحد، درس و استاد را اضافه، ویرایش و حذف می‌کند.
\item
  درخواست دانشجو برای افزودن استاد در پنل مدیر دیده و تأیید یا رد می‌شود.
\item
  حذف مورد دارای جزوه مجاز نیست و پیام خطا می‌دهد.
\end{itemize}

\end{document}
